\section{Votación ponderada: registro y resultados}
\label{sec:votacion}

\subsection{Datos y control de calidad}

El 27 de noviembre de 2026 se registraron tres formularios identificados como P06, P07 y P08, distintos de los cinco códigos de Kano. Cada formulario contiene una asignación de puntos enteros no negativos para los veinte requisitos (RF-01 a RF-12 y RNF-01 a RNF-08). Las tres sumas individuales son exactamente cien puntos, por lo que el conjunto agregado totaliza trescientos puntos.

Para obtener el puntaje total de un requisito se suman sus tres asignaciones. El rango 1 corresponde al mayor total. Los empates reciben el rango medio de las posiciones que ocupan; así, dos requisitos empatados en las posiciones 6 y 7 tienen rango 6,5. Se conservan los valores individuales en \path{datos/votos.csv} y el cálculo verificable en \path{datos/calculos/orden_votacion.csv}.

\subsection{Orden obtenido}

\begingroup
\small
\begin{longtable}{@{}p{0.13\linewidth}p{0.53\linewidth}rr@{}}
\caption{Votación ponderada y rango medio}\label{tab:resultados-votacion}\\
\toprule
\textbf{ID} & \textbf{Requisito} & \textbf{Puntos} & \textbf{Rango}\\
\midrule
\endfirsthead
\toprule
\textbf{ID} & \textbf{Requisito} & \textbf{Puntos} & \textbf{Rango}\\
\midrule
\endhead
\bottomrule
\endlastfoot
RF-04 & Matrícula & 31 & 1,0\\
RF-08 & Registro de calificación & 25 & 2,0\\
RF-05 & Control de cupos & 24 & 3,0\\
RF-01 & Registro de estudiante & 21 & 4,0\\
RF-10 & Consulta de notas & 19 & 5,0\\
RF-07 & Asignación docente & 18 & 6,5\\
RF-12 & Reporte de matrícula & 18 & 6,5\\
RF-09 & Corrección de calificación & 17 & 8,0\\
RF-02 & Autenticación & 16 & 9,0\\
RNF-03 & Privacidad de calificaciones & 15 & 10,0\\
RF-06 & Comprobante de matrícula & 13 & 11,0\\
RF-11 & Boletín académico & 12 & 12,5\\
RNF-02 & Disponibilidad & 12 & 12,5\\
RF-03 & Oferta de cursos & 11 & 14,5\\
RNF-01 & Rendimiento de matrícula & 11 & 14,5\\
RNF-04 & Facilidad de uso de matrícula & 9 & 16,5\\
RNF-05 & Compatibilidad móvil & 9 & 16,5\\
RNF-06 & Copia de seguridad & 8 & 18,0\\
RNF-07 & Recuperación del servicio & 6 & 19,0\\
RNF-08 & Trazabilidad de modificaciones & 5 & 20,0\\
\end{longtable}
\endgroup


RF-04 (matrícula) obtuvo 31 puntos; RF-08 (registro de calificaciones), 25; y RF-05 (control de cupos), 24. El primer requisito de calidad en la ordenación fue RNF-03 (privacidad de calificaciones), con 15 puntos. La posición relativa deriva únicamente de estas tres papeletas, no de una valoración de toda la comunidad educativa.
